\chapter{Enteros parte 2}

\section{Simplificación en congruencias}

\begin{proposicion}[{title={Regla 1: simplificar un factor en toda la congruencia (incluido el módulo)}}]
Sean $a,b,k,m\in\mathbb{Z}$ con $k\neq 0$ y $m\ge 1$. Entonces
\[
ak\equiv bk \pmod{mk} \iff a\equiv b \pmod m.
\]
Es decir, si los dos lados y el módulo comparten un divisor común $k$, se lo
puede cancelar de los tres.
\end{proposicion}

\begin{demostracion}
Por definición, $ak\equiv bk \pmod{mk}$ si y sólo si $mk\mid ak-bk=(a-b)k$,
es decir, existe $t\in\mathbb{Z}$ con $(a-b)k=mk\,t$. Como $k\neq 0$, esto es
equivalente a $a-b=mt$, o sea $m\mid a-b$, que es $a\equiv b \pmod m$.
\end{demostracion}

\begin{ejemplo}
Resolver $6x\equiv 8 \pmod{10}$.

Los tres números $6$, $8$ y $10$ son divisibles por $2$, así que podemos
``cancelar'' el $2$ dividiendo todo, incluido el módulo:
\[
\frac{6}{2}x\equiv \frac{8}{2} \pmod{\frac{10}{2}} \iff 3x\equiv 4 \pmod 5.
\]
Para terminar, como $3\cdot 2=6\equiv 1 \pmod 5$, multiplicamos por $2$ (que es
coprimo con $5$): $x\equiv 8\equiv 3 \pmod 5$. Chequeo: $6\cdot 3=18\equiv 8 \pmod{10}$. $\checkmark$
\end{ejemplo}

\begin{proposicion}[{title={Regla 2: simplificar un factor sólo a los lados (si es coprimo con el módulo)}}]
Sean $a,b,k,m\in\mathbb{Z}$ con $m\ge 1$. Si $(k:m)=1$, entonces
\[
ak\equiv bk \pmod m \iff a\equiv b \pmod m.
\]
Es decir, si queremos dividir sólo los números de la izquierda y la derecha por
$k$, el módulo no cambia \textbf{sólo si} $k$ y $m$ son coprimos.
\end{proposicion}

\begin{demostracion}
($\Leftarrow$) Es la compatibilidad del producto con la congruencia.

($\Rightarrow$) Si $ak\equiv bk \pmod m$, entonces $m\mid k(a-b)$. Como $(k:m)=1$,
por Bézout existen $x,y\in\mathbb{Z}$ con $xm+yk=1$. Multiplicando por $a-b$:
\[
a-b=xm(a-b)+y\,k(a-b).
\]
Como $m\mid xm(a-b)$ y $m\mid y\,k(a-b)$, resulta $m\mid a-b$.
\end{demostracion}

\begin{ejemplo}
Resolver $6x\equiv 12 \pmod 5$.

Como $(6:5)=1$, podemos dividir $6x$ y $12=6\cdot 2$ por $6$ sin tocar el
módulo $5$:
\[
6x\equiv 6\cdot 2 \pmod 5 \iff x\equiv 2 \pmod 5.
\]
\end{ejemplo}

\begin{nota}
Si $(k:m)\neq 1$ no se puede dividir sólo a los lados sin cambiar el módulo. Por
ejemplo $2\cdot 3\equiv 2\cdot 1 \pmod 4$ (pues $6\equiv 2$), pero
$3\not\equiv 1 \pmod 4$. Lo correcto en ese caso es la Regla 1, que también
divide el módulo: $3\equiv 1 \pmod 2$.
\end{nota}

\section{Teorema chino del resto}
\begin{proposicion}
Dado un sistema de ecuaciones de congruencia en $\mathbb{Z}$
\[
\begin{cases}
x\equiv c_1 \pmod{m_1}\\
x\equiv c_2 \pmod{m_2}\\
\quad\vdots\\
x\equiv c_n \pmod{m_n}
\end{cases}
\]
si $(m_i:m_j)=1$ para todo $i\neq j$ (los módulos son coprimos dos a dos),
entonces el sistema tiene una única solución $x_0$ tal que
\[
0\le x_0<m_1m_2\cdots m_n.
\]
Más aún, en $\mathbb{Z}$ hay infinitas soluciones, y el conjunto de soluciones es
\[
S=\{x\in\mathbb{Z} : x\equiv x_0 \pmod{m_1m_2\cdots m_n}\}.
\]
\end{proposicion}

\begin{demostracion}
\textbf{Existencia.} Por inducción en $n$; el caso $n=1$ es trivial. Veamos cómo
juntar dos ecuaciones $x\equiv c_1 \pmod{m_1}$ y $x\equiv c_2 \pmod{m_2}$ con
$(m_1:m_2)=1$. Las soluciones de la primera son $x=c_1+m_1t$, $t\in\mathbb{Z}$.
Reemplazando en la segunda queda $m_1t\equiv c_2-c_1 \pmod{m_2}$, que tiene
solución porque $(m_1:m_2)=1$ divide a todo, y sus soluciones forman una clase
$t\equiv t_0 \pmod{m_2}$. Entonces $x=c_1+m_1t_0+m_1m_2 j$, es decir, las dos
ecuaciones equivalen a una sola: $x\equiv x_1 \pmod{m_1m_2}$. Como $m_1m_2$ es
coprimo con $m_3,\dots,m_n$, se repite el argumento.

\textbf{Unicidad.} Si $x$ y $x'$ son soluciones, entonces $m_i\mid x-x'$ para
todo $i$. Como los $m_i$ son coprimos dos a dos, $m_1\cdots m_n\mid x-x'$. En
particular, dos soluciones entre $0$ y $m_1\cdots m_n-1$ son iguales.
\end{demostracion}

\section{Pequeño teorema de Fermat}

\begin{proposicion}[{title={Pequeño teorema de Fermat}}]
Sea $p$ un número primo y sea $a\in\mathbb{Z}$. Entonces:
\begin{enumerate}
    \item $a^p\equiv a \pmod p$.
    \item Si $p\nmid a$, entonces $a^{p-1}\equiv 1 \pmod p$.
\end{enumerate}
\end{proposicion}

\begin{demostracion}
\textbf{Prueba de (2).} Supongamos $p\nmid a$. Consideramos los $p-1$ números
\[
a,\ 2a,\ 3a,\ \dots,\ (p-1)a.
\]
\begin{itemize}
    \item Ninguno es múltiplo de $p$: si $p\mid ia$ con $1\le i\le p-1$, por el
    lema de Gauss $p\mid i$ o $p\mid a$, pero $0<i<p$ y $p\nmid a$.
    \item Sus restos módulo $p$ son todos distintos: si $ia\equiv ja \pmod p$
    con $1\le i,j\le p-1$, como $(a:p)=1$ podemos cancelar $a$ (Regla 2) y
    queda $i\equiv j \pmod p$, así que $i=j$ por ser ambos menores que $p$.
\end{itemize}
Entonces los restos de $a,2a,\dots,(p-1)a$ son $p-1$ valores distintos en
$\{1,\dots,p-1\}$, es decir, son exactamente $1,2,\dots,p-1$ en algún orden.
Multiplicándolos:
\[
a\cdot 2a\cdots (p-1)a\equiv 1\cdot 2\cdots (p-1) \pmod p
\iff a^{p-1}(p-1)!\equiv (p-1)! \pmod p.
\]
Como $p$ es primo y no divide a ninguno de los factores $1,\dots,p-1$, por el
corolario del lema de Gauss $p\nmid (p-1)!$, o sea $((p-1)!:p)=1$. Cancelando
$(p-1)!$ (Regla 2) queda
\[
a^{p-1}\equiv 1 \pmod p.
\]

\textbf{Prueba de (1).} Hay dos casos.
\begin{itemize}
    \item Si $p\nmid a$: por (2), $a^{p-1}\equiv 1 \pmod p$, y multiplicando por
    $a$ se obtiene $a^p\equiv a \pmod p$.
    \item Si $p\mid a$: entonces $a\equiv 0 \pmod p$ y $a^p\equiv 0 \pmod p$,
    así que $a^p\equiv a \pmod p$.
\end{itemize}
\end{demostracion}

\begin{proposicion}[{title={Corolario}}]
Sea $p$ un número primo y sea $a\in\mathbb{Z}$ con $p\nmid a$. Si $n\in\mathbb{N}_0$,
entonces
\[
a^n\equiv a^{r_{p-1}(n)} \pmod p,
\]
donde $r_{p-1}(n)$ es el resto de dividir $n$ por $p-1$.
\end{proposicion}

\begin{demostracion}
Por el algoritmo de la división, $n=q(p-1)+r$ con $r=r_{p-1}(n)$ y
$0\le r<p-1$. Como $p\nmid a$, por Fermat $a^{p-1}\equiv 1 \pmod p$, así que
\[
a^n=a^{q(p-1)+r}=\left(a^{p-1}\right)^q a^r\equiv 1^q\, a^r=a^r \pmod p.
\]
\end{demostracion}

\begin{nota}{}
El corolario requiere $p\nmid a$. Si $p\mid a$, no se puede reducir el exponente
de esta forma: por ejemplo, con $p=3$, $a=3$ y $n=2$ se tiene $r_2(2)=0$, pero
$3^0=1\not\equiv 0\equiv 3^2 \pmod 3$.
\end{nota}

\begin{ejemplo}
Calcular el resto de dividir $3^{100}$ por $7$.

Como $7$ es primo y $7\nmid 3$, por Fermat $3^6\equiv 1 \pmod 7$. Como
$100=6\cdot 16+4$, es decir $r_6(100)=4$, por el corolario
\[
3^{100}\equiv 3^{r_6(100)}=3^4=81\equiv 4 \pmod 7,
\]
pues $81=77+4$. El resto es $4$.
\end{ejemplo}